\chapter{Dissociative Identity Disorder}

\medskip
\noindent\textit{\textbf{Under review.} This chapter is an AI-assisted educational draft; it is not a substitute for professional medical advice.}

\section*{Definition and Overview}
Dissociative identity disorder (DID), formerly called multiple personality disorder, is a complex dissociative disorder involving distinct identity states and recurrent gaps in recall that are not explained by ordinary forgetfulness, substances, or another medical condition. Diagnosis requires careful specialist assessment.

\section*{Clinical Features}
People may experience memory gaps, depersonalisation, derealisation, identity confusion, distress, trauma-related symptoms, sleep problems, depression, anxiety, or self-harm risk. Experiences vary and should not be judged from media portrayals.

\section*{Differential Diagnosis}
Assessment distinguishes DID from post-traumatic stress disorder, borderline personality disorder, psychotic disorders, epilepsy, substance effects, sleep disorders, and neurological illness. Coexisting conditions are common.

\section*{When to Seek Medical Care}
Seek mental-health assessment for recurrent unexplained memory gaps or dissociation that causes distress or unsafe situations. Seek immediate help for suicidal thoughts, self-harm, inability to remain safe, or severe confusion.

\section*{Diagnosis and Management}
Diagnosis is clinical and should be made by a clinician experienced in dissociative disorders. Treatment is usually phased trauma-informed psychotherapy, stabilisation, safety planning, and treatment of coexisting conditions. Medication may target associated symptoms but does not remove identity states.

\section*{Self-Care}
Use grounding strategies taught by a clinician, maintain regular sleep and routines, avoid intoxicants, and involve trusted supports with consent. Emergency services should be contacted when safety cannot be maintained.

\section*{Expanded clinical context}
Multiple personality disorder is a mental-health topic. Interpretation should account for age, baseline health, medicines, comorbidities, goals, and access to care. This chapter supports discussion with a qualified clinician and does not establish a diagnosis.

\section*{Assessment and decision points}
Assessment should consider onset, duration, triggers, sleep, substance use, physical health, developmental context, social stressors, and immediate safety. Ask about self-harm or harm from others when appropriate. Document the main concern, time course, functional impact, relevant risks, and red flags. Use targeted investigations when their result is likely to change management.

\section*{Management and follow-up}
Management may include education, psychological therapy, practical support, treatment of coexisting conditions, medication when indicated, and an agreed follow-up or safety plan. Explain expected improvement, when reassessment is needed, and how follow-up can be accessed. Avoid starting, stopping, or sharing prescription medicines without professional advice.

\section*{Safety-netting}
Seek urgent help for immediate danger, suicidal intent, psychosis, severe confusion, inability to meet basic needs, or rapidly worsening symptoms.

\section*{Practical prevention}
Use plain-language information, supportive relationships, regular routines, and professional review when symptoms persist or impair daily life. Keep written instructions and seek review if the topic recurs, changes, or does not improve as expected.

\section*{Limitations}
This educational expansion should be checked against current local guidance and authoritative topic-specific sources.
\clearpage
